\begingroup
\raggedbottom
\setlength{\parskip}{6pt}
\section[Efficacité de formation stellaire]{Efficacité de formation stellaire}
\label{sec:sfe}

Les relations de Kennicutt--Schmidt décrivent les variations conjointes des densités de surface du gaz et du SFR. Suivant \citet{Leroy2008}, nous examinons comment l'efficacité apparente de formation stellaire diffère entre les réservoirs d'\HI\ étroit et large, et comment elle évolue dans le disque. L'efficacité de formation stellaire (SFE) est définie comme le SFR divisé par la masse de gaz.

Nous la calculons à partir des nouvelles cartes des composantes étroite et large, à la même résolution spatiale que les comparaisons individuelles de Kennicutt--Schmidt. Pour une région contenant $N_{\mathrm{reg}}$ pixels valides :
\begin{equation}
\label{eq:regional_sfe}
\mathrm{SFE}_p = 10^{-6}\,
\frac{\displaystyle\sum_{i=1}^{N_{\mathrm{reg}}}\Sigma_{\mathrm{SFR},i} A_i}
{\displaystyle\sum_{i=1}^{N_{\mathrm{reg}}}\Sigma_{p,i} A_i},
\end{equation}
\noindent où $A_i$ est l'aire du pixel et $p$ le traceur gazeux. Le facteur $10^{-6}$ convertit la densité de surface du SFR de kpc$^{-2}$ en pc$^{-2}$, de sorte que la SFE est exprimée en yr$^{-1}$. Pour comparer les traceurs sur une même aire, nous utilisons les valeurs des mêmes pixels dans les cartes de gaz et de SFR.

\begin{table}[tbp]
\centering
\caption{\textbf{M\'ediane de $\operatorname{\mathbf{log}}_{\mathbf{10}}\mathbf{SFE}$}, en yr$^{-1}$, calculée dans l'ouverture retenue de chaque galaxie avant le regroupement de l'échantillon. Les intervalles couvrent les 68\% centraux de 10\,000 rééchantillonnages bootstrap de galaxies entières, sans inclure les incertitudes de calibration et de décomposition. Les lignes d'\HI\ utilisent neuf naines et sept spirales. Les lignes incluant le gaz moléculaire utilisent les mêmes pixels couverts en CO dans quatre naines et sept spirales.}
\label{tab:sfe}
\begingroup

\renewcommand{\arraystretch}{1.3}
\begin{tabular}{lrrr}
\toprule
Gaz au dénominateur & Naines & Spirales & Ensemble \\
\midrule
\HI\ total & $-9.69^{+0.21}_{-0.10}$ & $-9.16^{+0.17}_{-0.13}$ & $-9.39^{+0.14}_{-0.29}$ \\
\HI\ étroit & $-9.29^{+0.21}_{-0.15}$ & $-8.66^{+0.05}_{-0.33}$ & $-9.08^{+0.19}_{-0.28}$ \\
\HI\ large & $-9.47^{+0.20}_{-0.06}$ & $-8.99^{+0.24}_{-0.00}$ & $-9.13^{+0.14}_{-0.27}$ \\
H$_2$ & $-8.45^{+0.19}_{-0.09}$ & $-9.10^{+0.13}_{-0.13}$ & $-8.96^{+0.39}_{-0.15}$ \\
H$_2$ + \HI\ étroit & $-9.02^{+0.13}_{-0.23}$ & $-9.25^{+0.02}_{-0.10}$ & $-9.23^{+0.17}_{-0.12}$ \\
H$_2$ + \HI\ large & $-8.99^{+0.04}_{-0.27}$ & $-9.35^{+0.09}_{-0.03}$ & $-9.26^{+0.25}_{-0.13}$ \\
H$_2$ + \HI\ total & $-9.22^{+0.06}_{-0.28}$ & $-9.43^{+0.01}_{-0.02}$ & $-9.42^{+0.20}_{-0.03}$ \\
\bottomrule
\end{tabular}
\endgroup
\end{table}

La figure~\ref{fig:adopted_sfe} montre la SFE en fonction du rayon pour les différentes composantes gazeuses. Chaque point représente la SFE médiane des galaxies contribuant à un intervalle radial de largeur $0.1r_{25}$. Dans les naines, le réservoir d'\HI\ étroit semble former les étoiles plus efficacement que le réservoir large dans l'ensemble du disque. Dans les spirales, la composante étroite présente une efficacité plus élevée dans le disque interne ($r<0.5r_{25}$), tandis que les deux efficacités se rapprochent à plus grand rayon. Dans les deux populations, la SFE apparente de l'\HI\ total et de chacune des composantes diminue vers l'extérieur, en accord avec la décroissance observée par \citet{Leroy2008} dans les régions dominées par l'\HI.

Dans les spirales, la SFE moléculaire varie relativement peu avec le rayon, alors que l'ajout de l'une ou l'autre composante atomique introduit une décroissance radiale et réduit les valeurs de SFE. Les profils H$_2$+\HI\ étroit et H$_2$+\HI\ large sont similaires.

Dans les régions les plus internes des spirales, ces deux combinaisons donnent des SFE proches de celle du H$_2$ seul, car le gaz des régions couvertes en CO y est principalement moléculaire. Ajouter une composante atomique modifie donc peu le résultat. À plus grand rayon, sa contribution augmente~: l'ajout d'\HI\ abaisse davantage la SFE et les profils combinés s'écartent du profil moléculaire, conformément à la diminution radiale de la fraction moléculaire observée par \citet{Leroy2008}. Dans les naines couvertes en CO, les deux combinaisons présentent déjà une SFE inférieure à celle du H$_2$ seul dans les régions centrales. Contrairement aux centres des spirales, ces régions contiennent une contribution atomique importante. Celle-ci augmentant vers l'extérieur, les deux profils combinés s'éloignent davantage du profil moléculaire.

Le tableau~\ref{tab:sfe} résume les SFE médianes des échantillons de naines et de spirales.

\begin{figure}[H]
\centering
\includegraphics[width=\linewidth]{images/ks_individual/adopted_sfe_profiles.pdf}
\caption[Efficacité de formation stellaire avec les nouvelles cartes cinématiques]{Profils radiaux de SFE des naines (en haut) et des spirales (en bas). À gauche : SFR observé divisé par l'\HI\ total, étroit ou large. À droite : H$_2$ et ses combinaisons avec les composantes atomiques ou l'\HI\ total, aux mêmes positions couvertes en CO. Chaque point représente la médiane des mesures des galaxies dans un intervalle de $0.1r_{25}$. Les bandes ombrées couvrent les 16e--84e centiles de la distribution entre galaxies pour l'\HI\ total (à gauche) et le H$_2$ (à droite) ; elles ne représentent pas l'incertitude sur la médiane.}
\label{fig:adopted_sfe}
\end{figure}
\FloatBarrier
\clearpage
\endgroup
